\chapter{1931: qué preguntó Hilbert y qué contestó Gödel}
\label{cap:hilbert}

Este libro cuenta una formalización, pero la formalización no se entiende sin la pregunta que la
originó. Y la pregunta no la hizo Gödel: la hizo Hilbert, treinta años antes, y era una pregunta
optimista.

\section{El programa de Hilbert}

A principios del siglo~XX las matemáticas habían pasado un susto. Las paradojas de la teoría de
conjuntos mostraron que razonar con confianza sobre objetos infinitos podía llevar a
contradicciones, y que la intuición no bastaba para detectarlas. La respuesta de Hilbert fue un
programa: si no podemos fiarnos de la intuición, escribamos las matemáticas como un sistema formal
—símbolos, axiomas, reglas— y demostremos \emph{sobre ese sistema}, con métodos que nadie discuta,
que es consistente.

\begin{leccion}[La palabra clave es «finitario»]
Hilbert no pedía cualquier demostración de consistencia. Pedía una \textbf{finitaria}: que no usara
ella misma el infinito que estaba tratando de justificar. Razonar sobre secuencias finitas de
símbolos, comprobar cosas caso a caso, contar. Nada más. Ésa era la condición, y es la que da
sentido a todo lo que viene después — incluida buena parte de las decisiones de este proyecto.
\end{leccion}

Para plantear siquiera la pregunta hay que separar dos cosas que en matemáticas ordinarias van
juntas. Por un lado, el sistema del que se habla: sus fórmulas, sus axiomas, su noción de
demostración. Lo llamaremos la \defterm{teoría objeto}, y sus fórmulas son, para nosotros,
\emph{datos}: cadenas de símbolos que se pueden manipular sin creerlas. Por otro, el lugar desde el
que hablamos de ella y demostramos cosas \emph{sobre} ella: la \defterm{metateoría}.

Esa separación no es un tecnicismo. Es la condición de posibilidad del programa entero: sin ella,
«demostrar que la aritmética es consistente» sería demostrarlo \emph{dentro} de la aritmética, y no
convencería a nadie.

\section{El finitismo era una estrategia, no una creencia}
\label{sec:finitismo-estrategia}

Conviene deshacer ahora un malentendido que se arrastra, porque cambia cómo se lee todo lo demás.
Hilbert \textbf{no era} un matemático constructivista que desconfiara del infinito. Era justo lo
contrario: había hecho carrera usando métodos fuertemente no constructivos y los defendía como un
territorio ganado, un derecho que las matemáticas no debían devolver.

El caso famoso es el teorema de la base, en teoría de invariantes: Hilbert demostró que existe una
base finita \emph{sin exhibir ninguna}, en un campo cuyo rey indiscutible, Paul Gordan, había
construido las suyas a mano. La reacción que se cita siempre es la de Gordan —«esto no es
matemática, esto es teología»—, y este libro, que se dedica a comprobar afirmaciones, tiene que
añadir en el mismo aliento que \textbf{esa cita está mal atestiguada}: aparece impresa por primera
vez en la necrológica que Max Noether escribió \emph{veinticinco años después} del episodio. Lo que
sí está documentado es lo contrario del mito: Gordan publicó en 1893 su propia versión del
argumento diciendo que la demostración de Hilbert era, en sustancia, correcta.

\begin{observacion}
Se cita esa frase para ilustrar la resistencia al método no constructivo, y funciona tan bien como
ilustración que casi nadie va a mirar de dónde sale. Es exactamente el modo de fallo que
\doc{PLAN-LIBRO.md}~§2.6 describe para los docstrings de este proyecto: una afirmación cómoda,
repetida hasta ser evidente, cuya fuente nadie ha vuelto a abrir. Se cita aquí porque es
instructiva, y con la salvedad delante.
\end{observacion}

Con eso en su sitio, la exigencia finitaria se lee bien. Hilbert no pedía que las
\emph{matemáticas} fueran finitarias — pedía que lo fuera la \textbf{demostración de
consistencia}, que es otra cosa. Se puede usar todo el infinito que se quiera \emph{dentro} del
sistema formal, con tal de que el argumento que certifica ese sistema desde fuera no lo use. Es una
autolimitación de la metateoría, y es estratégica: sólo un argumento que no presuponga lo que está
en duda puede convencer a quien lo duda.

\begin{leccion}[Y Gödel hizo lo mismo, por la misma razón]
La restricción de Gödel también era estratégica. Su argumento no usa el infinito ni necesita lógica
clásica más allá de lo imprescindible — no porque él dudara de ellos, sino porque un teorema que
demoliera el programa de Hilbert usando los métodos que el programa quería justificar no habría
demolido nada. Los dos se autolimitaron para poder hablarle al otro.

Esa es, tres capas más abajo, la misma razón por la que este proyecto anota debajo de cada teorema
de qué axiomas depende, y por la que el capítulo~\ref{cap:objeto-y-meta} se molesta en contar
cuánta lógica clásica gasta:
\textbf{importa con qué se demuestra, no sólo qué se demuestra.}
\end{leccion}

\section{La idea de Gödel}

Gödel publicó en 1930 un resultado que encajaba con el optimismo: la lógica de primer orden es
completa —todo lo que es verdad en todos los modelos es demostrable—. Un año después publicó el
que no encajaba.

La idea central es una jugada de una audacia difícil de exagerar. Si las fórmulas de la teoría
objeto son cadenas de símbolos, y las cadenas de símbolos se pueden numerar, entonces \textbf{la
teoría puede hablar de sus propias fórmulas} — no de fórmulas, sino de los números que las
codifican. A ese procedimiento se le llama \defterm{aritmetización} o \defterm{gödelización}, y el
número —o el término que lo denota— es el \defterm{código de Gödel} de la fórmula.

\begin{leccion}[Por qué eso lo cambia todo]
Una teoría aritmética habla de números. Si las fórmulas \emph{son} números, la teoría habla de
fórmulas sin saberlo. Y si además puede expresar «existe una demostración de la fórmula número
$x$», entonces habla de su propia demostrabilidad. La autorreferencia deja de ser un juego de
palabras y pasa a ser aritmética.
\end{leccion}

Con esa maquinaria, Gödel construye una sentencia $G$ que la teoría demuestra equivalente a «yo no
soy demostrable» — un \defterm{punto fijo} de la negación de la demostrabilidad — y prueba que si la
teoría es consistente, no puede demostrar $G$. El primer teorema. Y a continuación observa que el
argumento entero, siendo finitario, se puede formalizar dentro de la propia teoría: de donde
resulta que la teoría tampoco puede demostrar su propia consistencia. El segundo teorema.

\section{El «I» del título}

Y aquí está el detalle que da sentido a este libro.

El artículo de 1931 se titula \emph{Sobre proposiciones formalmente indecidibles de los Principia
Mathematica y sistemas afines}, \textbf{I}. Ese «I» no es decorativo: Gödel anunció una segunda
parte con la demostración detallada del segundo teorema, que en el artículo publicado sólo queda
esbozada. \textbf{La segunda parte nunca se escribió.}

Lo que faltaba no era la idea, que estaba clara, sino la comprobación penosa de que el predicado de
demostrabilidad se comporta como debe: que si algo es demostrable, la teoría demuestra que lo es;
que la demostrabilidad respeta el modus ponens; y que la teoría demuestra ella misma esa primera
propiedad. Esas tres son las \defterm{condiciones de derivabilidad}, y las suministraron
Hilbert y Bernays en 1939.

\begin{leccion}[La tesis de este libro, en una frase]
Las condiciones de derivabilidad son el artículo que Gödel no llegó a publicar. Y de las tres, la
tercera —la que dice que la teoría demuestra \emph{de sí misma} que reconoce lo que es cierto de
forma finita, la \defterm{$\Sigma_1$-completitud} en su versión \defterm{provable}— es la que este
proyecto lleva meses sin poder cerrar. Buena parte de la Parte~\ref{parte:noescrito} cuenta ese intento.
\end{leccion}

\section{Qué quedó del programa}

El programa de Hilbert, tal y como estaba formulado, no se puede completar. Pero decir sólo eso es
quedarse corto de una manera que este libro quiere evitar: lo que Gödel demostró no es que las
matemáticas sean poco fiables, sino que \textbf{la fiabilidad no se puede certificar desde dentro}.
Es un resultado sobre los límites de la autocertificación, no sobre la verdad matemática.

Y tiene una consecuencia práctica que hoy se ve mejor que en 1931: si un sistema no puede
certificarse a sí mismo, la alternativa es certificarlo desde otro sistema — explícito, más débil o
más fuerte, pero declarado. Que es exactamente lo que hace una formalización asistida por ordenador,
y lo que hace este proyecto.
